\documentclass[11pt]{article}

\def\chapitre{4}
\def\pagetitle{Espaces Vectoriels Normés.}

\input{setup.tex}



\begin{document}

\input{title.tex}

\thispagestyle{fancy}

L'idée du chapitre est de faire de l'analyse (définir la notion de limite) sur d'autres objets que les réels ou les complexes.

\section{Normes et distances.}

\begin{defi}{Norme. $\star$}{}
    Soit $E$ un $\K$-ev et $N:x\mapsto N(x)$ de $E$ dans $\R_+$. C'est une \bf{norme} sur $E$ ssi :
    \begin{enumerate}
        \item $\forall x \in E, ~ N(x)=0 \iff x = 0_E$ (\bf{Séparation})
        \item $\forall x \in E, ~ \forall \lambda \in \K, ~ N(\lambda x) = |\lambda|N(x)$ (\bf{Homogénéité})
        \item $\forall x,y \in E, ~ N(x+y) \leq N(x) + N(y)$ (\bf{Inégalité triangulaire})
    \end{enumerate}
\end{defi}

\vspace*{-0.8cm}

\begin{ex}{}{}
    La valeur absolue est une norme sur $\R$, le module est une \bf{norme} sur $\C$.
\end{ex}

\vspace*{-0.8cm}

\begin{ex}{}{}
    Montrer que toutes les normes sur $\R$ sont proportionnelles à la valeur absolue.
    \tcblower
    Soit $N$ une norme sur $\R$, alors $\forall \l \in \R, ~ N(\l)=|\l|N(1)$ par homogénéité.
\end{ex}

\vspace*{-0.8cm}

\begin{ex}{}{}
    Donner deux normes non proportionnelles sur $\C$.
    \tcblower
    Le module est une norme sur $\C$, et $N(a+ib)=|a|+|b|$ aussi; non proportionnelles.
\end{ex}

\begin{prop}{Généralisation de l'inégalité triangulaire.}{}
    Soit $E$ un $\K$-ev et $\|\cdot\|$ une norme sur $E$.
    \begin{enumerate}
        \item $\forall(x_1,...,x_n)\in E^n, ~ \|\sum_{i=1}^nx_i\|\leq\sum_{i=1}^n\|x_i\|$.
        \item $\forall x,y \in E, ~ |\|x\|-\|y\||\leq\|x-y\|$.
    \end{enumerate}
    \tcblower
    \boxed{1.} Par récurrence facile sur $n$.\\
    \boxed{2.} Soient $x,y\in E$.
    \begin{equation*}
        \begin{cases}
            x=x-y+y\\
            y=y-x+x
        \end{cases} \ra \begin{cases}
            \|x\|\leq\|x-y\|+\|y\|\\
            \|y\|\leq\|y-x\|+\|x\|
        \end{cases} \ra \begin{cases}
            \|x\|-\|y\| \leq \|x-y\|\\
            \|y\|-\|x\| \leq \|x-y\|
        \end{cases}
    \end{equation*}
    Donc $|\|x\|-\|y\|| \leq \|x-y\|$.
\end{prop}

\vspace*{-0.8cm}

\begin{defi}{Espace vectoriel normé.}{}
    Un \bf{espace vectoriel normé} (evn) est un espace vectoriel $E$ muni d'une norme $\|\.\|$, noté $(E,\|\cdot\|)$.
\end{defi}

\vspace*{-0.8cm}

\begin{defi}{Vecteur unitaire.}{}
    Un vecteur d'un evn est dit \bf{unitaire} ssi il est de norme 1.
\end{defi}

\vspace*{-0.8cm}

\begin{defi}{Distance.}{}
    Soit $E$ un ensemble. Alors $d:E\times E \to \R_+$ est une \bf{distance} sur $E$ ssi
    \begin{enumerate}
        \item $\forall (x,y)\in E^2, ~ d(x,y)=0 \iff x = y$. (séparation)
        \item $\forall (x,y)\in E^2, ~ d(x,y)=d(y,x)$. (symétrie)
        \item $\forall (x,y,z)\in E^3, ~ d(x,z)\leq d(x,y)+d(y,x)$. (inégalité triangulaire)
    \end{enumerate}
\end{defi}

\vspace*{-0.8cm}

\begin{prop}{Distance associée à une norme. $\star$}{}
    Soit $(E,\|\cdot\|)$ un evn, alors $d:(x,y)\mapsto \|x-y\|$ est une distance sur $E$.
    \tcblower
    \boxed{1.} $d(x,y)=0 \iff \|x-y\|=0 \iff x-y = 0 \iff x = y$.\\
    \boxed{2.} $d(x,y)=\|x-y\|=|-1|\|y-x\|=\|y-x\|=d(y,x)$.\\
    \boxed{3.} $d(x,z)=\|x-z\|=\|x-y+y-z\|\leq \|x-y\| + \|y-z\| = d(x,y) + d(y,z)$.
\end{prop}

\vspace*{-0.8cm}

\begin{rappel}{}{}
    Un produit scalaire sur $E$ est une forme bilinéaire, symétrique, définie, positive.
\end{rappel}

\vspace*{-0.8cm}

\begin{prop}{Norme associée à un produit scalaire. $\star$}{}
    Soit $(E,\Lg\.,\.\Rg)$ un espace préhilbertien réel. On pose $\|x\|=\sqrt{\langle x,x \rangle}$, alors $\|\cdot\|$ définit une norme sur $E$.\\
    \bf{Remarque.} Si $\|\cdot\|$ est définie avec une racine carrée, on passe par le produit scalaire.
\end{prop}

\vspace*{-0.8cm}

\begin{exi}{$\K^n$. $\star$}{}
    On définit trois normes (à connaître) sur $\K^n$ :
    \begin{enumerate}[itemsep=0.1cm]
        \item $\forall (x_1,...,x_n) \in E^n, ~ \|(x_1,...,x_n)\|_1=\sum_{i=1}^n|x_i|$.
        \item $\forall (x_1,...,x_n) \in E^n, ~ \|(x_1,...,x_n)\|_2=\sqrt{\sum_{i=1}^n|x_i|^2}$.
        \item $\forall (x_1,...,x_n) \in E^n, ~ \|(x_1,...,x_n)\|_\infty=\max(|x_i|)$.
    \end{enumerate}
\end{exi}

\vspace*{-0.4cm}

\begin{lemme}{Homogénéité du sup. $\star$}{}
    Soit $A$ une partie non-vide et majorée de $\R$, et $\l\in\R_+^*$. Alors :
    \begin{equation*}
        \sup(\l A) = \l\sup(A).
    \end{equation*}
    \tcblower
    \boxed{\leq} Soit $a \in \l A ~ : ~ \exists x \in A, ~ a = \l x$, et $x\leq \sup(A)$ puisque $x\in A$.\\
    Donc $\l x \leq \l \sup(A)$ puisque $\l>0$, puis $a \leq \l \sup(A)$. Par passage au sup, $\sup(\l A) \leq \l \sup(A)$.\\
    \boxed{\geq} Soit $x\in A$, alors $\l x \in \l A$ donc $\l x \leq \sup(\l A)$ puis $x \leq \frac{1}{\l}\sup(\l A)$.\\
    Par passage au sup, $\sup(A) \leq \frac{1}{\l}\sup(\l A)$ donc $\l\sup(A) \leq \sup(\l A)$.\\
    \bf{Conclusion.} Par double-inégalité, $\sup(\l A) = \l\sup(A)$. 
\end{lemme}

\vspace*{-0.8cm}

\begin{exi}{Fonctions bornées. $\star$}{}
    On note $\m{B}(A,\R)$ l'ensemble des fonctions de $A\subset \R$ bornées sur $A$ non vide.\\
    Montrer que $\|f\|=\sup_{x\in A}|f(x)|$ définit une norme sur $\m{B}(A,\R)$.
    \tcblower 
    $\hra$ $\{|f(x)|, ~ x \in A\}$ est un sous-ensemble non-vide et majoré de $\R$ : il admet une borne supérieure.\\
    \bf{Séparation.} $\|f\|=0\ra\forall x \in A, ~ 0 \leq |f(x)| \leq \|f\|=0$, donc $\forall x \in A, f(x)=0$ puis $f=\zero_A$.\\
    \bf{Homogénéité.} Pour $\l\in\R$, $\|\l f\| = \sup_{x\in A}|\l f(x)| = |\l|\sup_{x\in A}|f(x)| = |\l|\|f\|$ par le lemme précédent.\\
    \bf{I.T.} $|(f+g)(x)|=|f(x)+g(x)|\leq|f(x)|+|g(x)|\leq \|f\|+\|g\|$ donc $\|f+g\|\leq\|f\|+\|g\|$    
\end{exi}

\vspace*{-0.8cm}

\begin{exi}{Fonctions continues. $\star$}{}
    On définit trois normes sur $\m{C}^0([a,b],\K)$, avec $a<b$ réels :
    \begin{enumerate}
        \item $\forall f \in \m{C}^0([a,b],\K), ~ \|f\|_1=\int_a^b|f(t)|\dt$.
        \item $\forall f \in \m{C}^0([a,b],\K), ~ \|f\|_2=\sqrt{\int_a^b|f(t)|^2\dt}$.
        \item $\forall f \in \m{C}^0([a,b],\K), ~ \|f\|_\infty=\sup\{|f(x)|, ~ x\in[a,b]\}$.
    \end{enumerate}
\end{exi}

\vspace*{-0.8cm}

\begin{ex}{}{}
    Pour $P\in\R[X]$, on définit $\|P\|=\int_0^1|P(t)|\dt$, une norme sur $\R[X]$.
    \tcblower
    \bf{Séparation.} $\|P\|=0 \ra \forall t \in [0,1], ~ P(t)=0$ car $|P|$ positive et continue, donc $P=0$ par rigidité.\\
    \bf{Homogénéité.} $\|\l P\|=\int_0^1|\l P(t)|\dt = |\l|\int_0^1|P(t)|\dt=|\l|\|P\|$.\\
    \bf{I.T.} $|(P+Q)(t)|\leq|P(t)|+|Q(t)|$, puis par croissance et linéarité : $\|P+Q\|\leq\|P\|+\|Q\|$.
\end{ex}

\vspace*{-0.8cm}

\begin{ex}{}{}
    Soient $a_1<...<a_n\in\R$. Pour $P\in\R[X]$, on definit $\|P\|=\sum_{i=0}^n|P(a_i)|$.
\end{ex}

\vspace*{-0.8cm}

\begin{ex}{}{}
    On définit trois normes sur $M_n(\R)$ :
    \begin{enumerate}
        \item $\forall A \in M_n(\R), ~ \|A\|_1=\sum_{1\leq i,j\leq n}|a_{i,j}|$.
        \item $\forall A \in M_n(\R), ~ \|A\|_2=\sqrt{\sum_{1\leq i,j \leq n} a_{i,j}^2}$.
        \item $\forall A \in M_n(\R), ~ \|A\|_{\infty}=\sup\{|a_{i,j}|, ~ 1 \leq i,j \leq n\}$.
    \end{enumerate}
\end{ex}

\vspace*{-0.8cm}

\begin{prop}{Normes sur les espaces produits.}{}
    Soit $(E_1,\|\cdot\|_1)$, ..., $(E_n,\|\cdot\|_n)$ des evn.\\
    On définit $\|\cdot\|_{\infty}$ sur $E=E_1\times...\times E_n$ par $\|(x_1,...,x_n)\|_{\infty}=\max(\|x_i\|_i, ~ 1\leq i \leq n)$.\\
    Alors $\|\cdot\|_{\infty}$ est une norme sur $E$, et $(E,\|,\|_{\infty})$ est l'evn produit.
    \tcblower
    \bf{Séparation.} $\|x\|_{\infty}=0\ra\forall i \in \lb1,n\rb, ~ \|x_i\|_i=0 $ donc $x_i=0_{E_i}$ donc $x=0_E$.\\
    \bf{Homogénéité.} $\|\l x\|_\infty = \max(\|\l x_i\|_i, ~ 1 \leq i \leq n) = |\l|\max(\|x_i\|_i, ~ 1 \leq i \leq n)=\l\|x\|_{\infty}$.\\
    \bf{I.T.} $\|x_i+y_i\|\leq \|x\|_{\infty} + \|y\|_\infty$ donc $\|x+y\|_\infty\leq\|x\|_{\infty}+\|y\|_{\infty}$.
\end{prop}

\vspace*{-0.8cm}

\section{Boules ouvertes, fermées, sphères.}

\begin{defi}{Boule ouverte, fermée et sphère.$\star$}{}
    Soit $(E,\|\cdot\|)$ un evn, $a\in E$ et $r\in\R_+$.
    \begin{enumerate}
        \item La \bf{boule ouverte} de centre $a$ et de rayon $r$ est l'ensemble $B(a,r)=\{x\in E, ~ \|x-a\| < r\}$.
        \item La \bf{boule fermée} de centre $a$ et de rayon $r$ est l'ensemble $B'(a,r)=\{x\in E, ~ \|x-a\|\leq r\}$.
        \item La \bf{sphère} de centre $a$ et de rayon $r$ est l'ensemble $S(a,r)=\{x\in E, ~ \|x-a\|=r\}$.
    \end{enumerate}
    \bf{Remarque.} $B(a,r)\sqcup S(a,r) = B'(a,r)$ et $r_1<r_2 \ra B(a,r_1)\subset B(a,r_2)$.\\
    \bf{Remarque.} Dans $\R$, les boules sont des intervalles ouverts ou fermés.
\end{defi}

\vspace*{-0.8cm}

\begin{nota}{}{}
    On appelle $B(0,1)$ la boule unité ouverte, $B'(0,1)$ la boule unité fermée, et $S(0,1)$ la sphère unité.
\end{nota}

\vspace*{-0.8cm}

\begin{defi}{}{}
    Soit $(E,\|\cdot\|)$ et $A\subset E$. On dit que $A$ est bornée ssi $\exists M > 0, ~ \forall x \in A, ~ \|x\|\leq M$.\\
    \bf{Remarque.} $A\subset B'(0,M)$.
\end{defi}

\vspace*{-0.8cm}

\begin{ex}{}{}
    Montrer que $A$ bornée $\iff \exists a \in E, ~ \exists r \geq 0, ~ A\subset B(a,r)$.
    \tcblower
    \boxed{\ra} Évident en posant $a=0$ et $r=M$.\\
    \boxed{\la} Soit $x\in A$, on a $\|x\|\leq \|a\| + \|x - a\| \leq \|a\| + r$. Donc $A$ est bornée en posant $M=\|a\|+r$. 
\end{ex}

\vspace*{-0.8cm}

\begin{defi}{Fonction bornée.}{}
    Soient $A$ un ensemble, $F$ un evn et $f:A\to F$. On dit que $f$ est \bf{bornée} ssi $f(A)$ est bornée dans $F$.
\end{defi}

\vspace*{-0.8cm}

\begin{defi}{Partie convexe.}{}
    Soit $E$ un $\K$-ev et $A\subset E$.\\
    On dit que $A$ est \bf{convexe} ssi $\forall (a,b) \in A^2, \forall \l \in [0,1], ~ \l a + (1-\l)b \in A$.
\end{defi}

\vspace*{-0.8cm}

\begin{rappel}{}{}
    Les parties convexes de $\R$ sont exactement les intervalles de $\R$.
\end{rappel}

\vspace*{-0.4cm}

\begin{prop}{}{}
    Soit $E$ un $\K$-ev et $A\subset E$ convexe. Une partie convexe est stable par passage aux barycentres :
    \begin{equation*}
        \forall n \in \N^*, ~ \forall (a_1,...,a_n) \in A^n, ~ \forall (\l_1,...,\l_n) \in \R_+^n, ~ \sum_{i=1}^n \l_i = 1 \ra \sum_{i=1}^n\l_ia_i\in A.
    \end{equation*}
    \tcblower
    Par récurrence sur $n\in\N^*$.\\
    \bf{Initialisation.} Pour $n\leq2$, c'est la définition de la convexité.\\
    \bf{Hérédité.} Soit $n\in\N^*$ tel que la propriété soit vraie, et $(a_1,...,a_{n+1})\in A^{n+1}$.\\
    Soit $(\l_1,...,\l_{n+1})\in\R_{+}^{n+1}$ une famille de réels positifs sommant à 1.\\
    $\hra$ Si $\l_{n+1}=1$, alors $\forall i \in \lb1,n\rb, ~ \l_i = 0$ puis $\sum_{i=1}^{n+1}\l_ia_i=a_{n+1}\in A$.\\
    $\hra$ Sinon, $\l_{n+1}<1$ puis :
    \begin{equation*}
        \sum_{i=1}^{n+1}\l_ia_i = (1-\l_{n+1})\sum_{i=1}^n \frac{\l_i}{1-\l_{n+1}}a_i + \l_{n+1}a_{n+1}.
    \end{equation*}
    Or $\forall i \in \lb1,n\rb, ~ \frac{\l_i}{1-\l_{n+1}} \geq 0$ et $\sum_{i=1}^n\frac{\l_i}{1-\l_{n+1}}=1$, donc par hypothèse de récurrence :
    \begin{equation*}
        \sum_{i=1}^n \frac{\l_i}{1-\l_{n+1}} a_i \in A ~\nt{puis}~ (1-\l_{n+1})\sum_{i=1}^n\frac{\l_i}{1-\l_{n+1}}a_i+\l_{n+1}a_{n+1}\in A
    \end{equation*}
    car $A$ est convexe.
\end{prop}

\vspace*{-0.8cm}

\begin{prop}{$\star$}{}
    Soit $(E,\|\cdot\|)$ un evn.
    \begin{enumerate}
        \item Les boules fermées sont convexes.
        \item Les boules ouvertes sont convexes.
    \end{enumerate}
    \tcblower
    Soient $a\in E$, $r\geq0$, $\l\in[0,1]$ et $x,y\in B(a,r)$.
    \begin{equation*}
        \|\l x + (1-\l)y - a\| = \|\l(x-a) + (1-\l)(y-a)\| \leq \l\|x-a\| + (1-\l)\|y-a\| < \l r + (1-\l) r = r.
    \end{equation*}
    Donc $\l x + (1-\l)y \in B(a,r)$ : c'est un convexe.
\end{prop}

\vspace*{-0.8cm}

\begin{ex}{}{}
    Vérifier que $S(a,r)$ pour $r>0$ n'est pas convexe.
    \tcblower
    Soit $x\in E$ non-nul.\\
    $\hra$ $y_1 = a - \frac{r}{\|x\|}x\in S(a,r)$ et $y_2 = a + \frac{r}{\|x\|}x\in S(a,r)$.\\
    $\hra$ $\frac{1}{2}y_1 + (1-\frac{1}{2})y_2 = a \notin S(a,r)$, donc $S(a,r)$ n'est pas convexe.
\end{ex}

\section{Suites dans un EVN.}

\begin{defi}{Suite bornée. $\star$}{}
    Soit $(E,\|\cdot\|)$ un evn et $(u_n)_{n\in\N}$ une suite de $E$.
    \begin{center}
        $(u_n)$ est \bf{bornée} ssi $\exists M \geq 0, ~ \forall n \in \N, ~ \|u_n\| \leq M$.
    \end{center}
\end{defi}

\vspace*{-0.8cm}

\begin{defi}{Convergence d'une suite. $\star$}{}
    Soit $(E,\|\cdot\|)$ un evn, $(u_n)_{n\in\N}$ une suite de $E$, et $l\in E$.
    \begin{center}
        $(u_n)_{n\in\N}$ \bf{converge} vers $l$ ssi $\lim \|u_n - l\| = 0$ ssi $\forall \e > 0, ~ \exists n_0 \in \N, ~ \forall n \geq n_0, ~ \| u_n - l \| \leq \e$.
    \end{center}
\end{defi}

\vspace*{-0.8cm}

\begin{ex}{}{}
    Soit $A\in M_n(\R)$, pour $k\in\N$, posons $A_k=A+\frac{1}{k+1}I_n$. Alors $\|A_k-A\|=\frac{1}{k+1}\|I_n\|\to 0$ puis $\lim A_k = A$.
\end{ex}

\vspace*{-0.8cm}

\begin{prop}{Unicité de la limite.}{}
    La limite d'une suite, si elle existe, est unique.
    \tcblower
    Soit $(u_n)$ une suite, convergente vers $l$ et $l'$. On suppose $l\neq l'$ par l'absurde.\\
    Soit $\e=\frac{1}{2}\|l-l'\| > 0$.\\
    $\hra$ $\exists n_0 \in \N, ~ \forall n \geq n_0, ~ \|u_n-l\|<\e$.\\
    $\hra$ $\exists n_1 \in \N, ~ \forall n \geq n_1, ~ \|u_n-l'\|<\e$.\\
    Soit $N = \max(n_0,n_1)$, alors $\forall n \geq N, ~ \|u_n-l\| < \e \et \|u_n-l'\|<\e$.\\
    $\hra$ $\|l-l'\| = \|(u_n - l) - (u_n - l')\| \leq \|u_n - l\| + \|u_n - l'\| < 2\e = \|l-l'\|$, absurde.
\end{prop}

\vspace*{-0.4cm}

\begin{thm}{Caractère borné d'une suite convergente. $\star$}{}
    Toute suite convergente est bornée.
    \tcblower
    Soit $(u_n)$ une suite de $(E,\|\cdot\|)$, supposée convergente vers $l\in E$.\\
    $\hra$ Par convergence de $(u_n)$ : $\exists n_0 \in \N, ~ \forall n \geq n_0, ~ \|u_n - l\| < 1$.\\
    Posons alors $r=\max(1,\|u_1-l\|,...,\|u_{n_0}-l\|)$.\\
    $\hra$ Pour $n\geq n_0$, $\|u_n - l\| \leq 1\leq r$ donc $u_n\in B(l,r)$.\\
    $\hra$ Pour $n\leq n_0, ~ \|u_n - l\|\leq r$ par définition de $r$, donc $u_n\in B(l,r)$.\\
    Donc $\forall n \in \N, ~ u_n \in B(l,r)$ : $(u_n)$ est bornée.
\end{thm}

\vspace*{-0.4cm}

\begin{prop}{$\star$}{}
    Soit $(E,\|\cdot\|)$ un evn, $(u_n)_{n\in\N}$, $(v_n)_{n\in\N}$ deux suites de $E$ convergentes, et $\l\in\K$.\\
    Alors $(u_n + \l v_n)_{n\in\N}$ est convergente et $\lim(u_n + \l v_n)=\lim u_n + \l \lim v_n$.
    \tcblower
    Supposons $\lim u_n = l$ et $\lim v_n = l'$.
    \begin{equation*}
        \|(u_n + \l v_n) - (l + \l l')\| = \|(u_n - l) + \l(v_n - l')\| \leq \|u_n - l\| + |\l|\|v_n - l'\| \to 0
    \end{equation*}
    Par encadrement, on a le résultat.
\end{prop}

\vspace*{-0.4cm}

\begin{defi}{Norme d'algèbre.}{}
    Soit $(A,+,\times,\cdot)$ une $\K$-algèbre, et $\|\cdot\|$ une norme sur $(A,+,\cdot)$.\\
    On dit que $\|\cdot\|$ est une \bf{norme d'algèbre} ssi $\forall (x,y)\in A, ~ \|x\times y\| \leq \|x\|\cdot\|y\|$.\\
    Une algèbre est dite normée si elle est munie d'une norme d'algèbre.
\end{defi}

\vspace*{-0.7cm}

\begin{prop}{}{}
    Soit $(A,+,\times,\cdot)$ une algèbre normée et $(u_n)_{n\in\N}$, $(v_n)_{n\in\N}$ convergentes.\\
    Alors $(u_n\times v_n)_{n\in\N}$ est convergente et $\lim (u_n\times v_n) = \lim u_n \times \lim v_n$.
    \tcblower
    Notons $l=\lim u_n$ et $l'=\lim v_n$. On a $(v_n)$ convergente donc bornée par $M$.
    \begin{equation*}
        \|u_nv_n - ll'\| = \|(u_n - l)v_n + l(v_n - l')\| \leq \|u_n - l\|\|v_n\| + \|l\|\|v_n - l'\| \leq \|u_n-l\|M + \|l\|\|v_n - l'\| \to 0 
    \end{equation*}
    Par encadrement, on a le résultat.
\end{prop}

\vspace*{-0.7cm}

\begin{prop}{Convergence dans un espace produit. $\star$}{}
    Soient $(E_1,\|\cdot\|_1),...,(E_N,\|\cdot\|_N)$ des evn et $E=E_1\times ... \times E_N$ l'evn produit.\\
    Soit $(u_n)_{n\in\N}$ une suite de $E$, $\forall n \in \N, ~ u_n = (u_n^1, ..., u_n^N)$.\\
    Soit $l=(l_1,...,l_N)\in E$. On a $u_n \to l \iff \forall i \in \lb1,N\rb, ~ (u_n^i)\to l_i$ dans $(E_i,\|\cdot\|_i)$.
    \tcblower
    \boxed{\ra} Soit $i\in\lb1,N\rb$. Par définition, $\|x\|_{\infty}=\max(\|x_i\|_i, 1 \leq i \leq N)$. Donc $\forall n \in \N, ~ \|u_n^i-l_i\|_i\leq \|u_n-l\|_{\infty}$.\\
    Par encadrement, $\|u_n^i-l_i\|\to0$.\\
    \boxed{\la} Supposons $\forall i \in \lb 1, N \rb, ~ \|u_n^i-l_i\|_i\to0$.\\
    Soit $\e>0, ~ \forall i \in \lb1,n\rb, ~ \exists p_i \in \N, ~ \forall n \geq p_i, ~ \|u_n^i - l_i\| \leq \e$.\\
    On prend alors $p=\max p_i$, puis pour $n\geq p$,  $\forall i \in \lb1,N\rb, ~ \|u_n^i-l_i\|\leq \e$ donc $\|u_n^i-l\|_{\infty}\leq\e$. 
\end{prop}

\vspace*{-0.7cm}

\begin{prop}{}{}
    Soient $E$ un $\K$-evdf de dimension $N$, et $\B=(e_1,...,e_n)$ une base de $E$.\\
    Soient $(u_n)_{n\in\N}\in E^\N$, et $l\in E$.
    \begin{equation*}
        \forall n \in \N, ~ u_n = \sum_{i=1}^Nu_n^ie_i ~\et~ l=\sum_{i=1}^Nl^ie_i 
    \end{equation*}
    On munit $E$ de la norme $\|x\|_\infty=\max(|x_i|, ~ 1\leq i \leq N)$, où $x=\sum x_i e_i$.\\
    Alors $u_n\to l$ dans $(E,\|\cdot\|_\infty) \iff \forall i \in \lb1,N\rb, ~ u_n^i \to l^i$.
\end{prop}

\vspace*{-0.7cm}

\begin{ex}{}{}
    Pour tout $n\in\N^*$, on pose la matrice :
    \begin{equation*}
        A_n = \begin{pmatrix}
            \frac{1}{n}+2 & -1 & e^{-n} \\ \ln\left( 2 + \frac{1}{n} \right) & 0 & \frac{n+1}{2n+3}
        \end{pmatrix}
    \end{equation*}
    $(A_n)_{n\in\N^*}$ est une suite de $M_{2,3}(\R)$ normé avec la norme produit associée à la base canonique, et on a :
    \begin{equation*}
        \lim_{n\to+\infty} A_n = \begin{pmatrix}
            2 & -1 & 0 \\ \ln 2 & 0 & \frac{1}{2}
        \end{pmatrix}
    \end{equation*}
\end{ex}

\begin{defi}{Suites extraites.}{}
    Soit $E$ un ensemble et $(u_n)_{n\in\N}\in E^\N$, alors $(v_n)_{n\in\N}$ est une \bf{suite extraite} de $(u_n)_{n\in\N}$ si et seulement si :
    \begin{equation*}
        \exists \phi : \N \to \N \nt{ strictement croissante telle que } \forall n \in \N, ~ v_n = u_{\phi(n)}.
    \end{equation*}
\end{defi}

\vspace*{-0.7cm}

\begin{lemme}{}{}
    Soit $\phi:\N\to\N$ strictement croissante, alors $\forall n \in \N, ~ \phi(n)\geq n$.
    \tcblower
    Par récurrence sur $n$, on note $\P_n$ : << $\phi(n)\geq n$ >>.\\
    \bf{Initialisation.} $\phi(0)\geq0$, donc $P_0$.\\
    \bf{Hérédité.} Soit $n\in\N$ tel que $\phi(n)\geq n$, alors $\phi(n+1)>\phi(n)\geq n$, donc $\phi(n+1)\geq n+1$.\\
    On conclut par récurrence.
\end{lemme}

\vspace*{-0.7cm}

\begin{defi}{Double extraction.}{}
    Soient $(u_n)_{n\in\N}$ et $(u_{\phi(n)})_{n\in\N}$, avec $\phi:\N\to\N$ strictement croissante.\\
    Une suite extraite de $(u_{\phi(n)})_{n\in\N}$ est de la forme $(u_{(\phi\circ\psi)(n)})_{n\in\N}$, avec $\psi:\N\to\N$ strictement croissante.
\end{defi}

\vspace*{-0.7cm}

\begin{prop}{}{}
    Soit $(E,\|\cdot\|)$ un evn et $(u_n)_{n\in\N}\in E^\N$.\\
    Si $(u_n)_{n\in\N}$ converge, alors toutes ses suites extraites convergent vers la même limite.\\
    Si $(u_n)_{n\in\N}$ possède deux suites extraites convergentes vers des limites distinctes, alors $(u_n)_{n\in\N}$ diverge.
    \tcblower
    Supposons $(u_n)_{n\in\N}$ convergente vers $l\in E$. Soit $\phi:\N\to\N$ strictement croissante.\\
    Pour $\e>0$, $\exists n_0 \in \N, ~ \forall n \geq n_0, ~ \|u_n-l\| < \e$, puis $\phi(n_0)\geq n_0$ donc $\|u_{\phi(n_0)}-l\|\leq \|u_{n_0} - l\| < \e$ donc $u_{\phi(n)}\to l$.
\end{prop}

\vspace*{-0.7cm}

\begin{prop}{}{}
    Soit $(E,\|\cdot \|)$ un evn et $(u_n)_{n\in\N}\in E^\N$. Soit $l\in\R$.
    \begin{equation*}
        \left[ u_{2n} \to l ~\et~ u_{2n+1} \to l \right] \ra u_n \to l.
    \end{equation*}
    \bf{Remarque.} Ce théorème est valable pour n'importe quelle partition finie des indices.
    \tcblower
    Soit $\e>0$.\\
    $\hra$ $u_{2n}\to l$ donc $\exists n_0 \in \N, ~ \forall n \geq n_0, ~ \|u_{2n}-l\|\leq \e$.\\
    $\hra$ $u_{2n+1}\to l$ donc $\exists n_1 \in \N, ~ \forall n \geq n_1, ~ \|u_{2n+1}-l\|\leq \e$.\\
    Soit $N=\max(n_0,n_1)$, alors $\forall n \geq N, ~ \|u_n - l\| \leq \e$.
\end{prop}

\vspace*{-0.7cm}

\begin{ex}{}{}
    Soient $(E,\|\cdot\|)$ un evn et $(u_n)_{n\in\N}\in E^\N$ tels que $(u_{2n})_{n\in\N}, (u_{2n+1})_{n\in\N}, (u_{3n})_{n\in\N}$ convergent.\\
    Montrer que $(u_n)_{n\in\N}$ converge.
    \tcblower
    On note $l=\lim u_{2n}$, $l'=\lim u_{2n+1}$ et $l''=\lim u_{3n}$.\\
    On a $(u_{6n})$ suite extraite de $(u_{2n})$ et de $(u_{3n})$, donc $\lim u_{6n}=l=l''$.\\
    On a $(u_{6n+3})$ suite extraite de $(u_{2n+1})$ et de $(u_{3n})$, donc $\lim u_{6n+3}=l'=l''$.\\
    Alors $l=l'=l''$, donc $\lim u_n = l$.
\end{ex}

\vspace*{-0.7cm}

\begin{ex}{}{}
    Soit $(u_n)_{n\in\N}$ une suite réelle non majorée positive.
    \begin{enumerate}
        \item Donner un exemple d'une telle suite qui ne tend pas vers $+\infty$.
        \item Montrer que l'on peut extraire de $(u_n)_{n\in\N}$ une suite qui tend vers $+\infty$.
    \end{enumerate}
    \tcblower
    \boxed{1.} Par exemple, $u_{2n}=666^n$ et $u_{2n+1}=1$.\\
    \boxed{2.} $\forall M \in \R_+, ~ \exists n \in \N, ~ u_n \geq M$.\\
    Soit $M\geq0$, $\exists n_0 \in \N, ~ u_{n_0}>M$, puis avec $M=u_{n_0}$, $\exists n_1, ~ u_{n_1}\geq u_{n_0}$...\\
    On a donc une infinité de majorants pour tout $M\geq0$. On construit $\phi$.\\
    --- $\exists n_0 \in \N, ~ u_{n_0}>0$, on pose $\phi(0)=n_0$.\\
    --- $\exists n_1 \in \N, ~ u_{n_1}>1$, on pose $\phi(1)=n_1$, on peut choisir $n_1>n_0$ car infinité d'entiers.\\
    --- $\exists n_p \in \N, ~ u_{n_p}>p$, on pose $\phi(p)=n_p$, on peut choisir $n_p>n_{p-1}$.
\end{ex}

\vspace*{-0.7cm}

\begin{defi}{Valeur d'adhérence.}{}
    Soient $(E,\|\cdot\|)$ un evn, $(u_n)_{n\in\N}\in E^\N$ et $l\in E$.\\
    On dit que $l$ est une \bf{valeur d'adhérence} de $(u_n)_{n\in\N}$ ssi $l$ est la limite d'une suite extraite de $(u_n)_{n\in\N}$.\\
    On notera $\Val((u_n)_{n\in\N})$ l'ensemble des valeurs d'adhérence de la suite $(u_n)_{n\in\N}$.
\end{defi}

\vspace*{-0.7cm}

\begin{ex}{}{}
    Soit $(u_n)_{n\in\N}$ telle que :
    \begin{itemize}
        \item $\forall n \in \N, u_n = n$, alors $\Val((u_n)_{n\in\N})=\0$.
        \item $\forall n \in \N, ~ \begin{cases} u_{2n} = 0 \\ u_{2n+1} = n \end{cases} \ra \Val((u_n)_{n\in\N}) = \{0\}$.
        \item $\forall n \in \N, ~ u_n=(-1)^n$, alors $\Val((u_n)_{n\in\N})=\{-1,1\}$.
        \item $\exists p \in \N, ~ \forall n \in \N, ~ u_{n+p}=u_n$, alors $\Val((u_n)_{n\in\N})=\{u_0,...,u_{p-1}\}$.
    \end{itemize}
\end{ex}

\vspace*{-0.7cm}

\begin{prop}{}{}
    Soit $(u_n)_{n\in\N}\in E^\N$ convergente de limite $l$, alors $\Val((u_n)_{n\in\N})=\{l\}$.\\
    La réciproque est fausse à priori : $u_{2n}=0$; $u_{2n+1}=n$, $0$ est la seule valeur d'adhérence.
\end{prop}

\vspace*{-0.7cm}

\section{Comparaison de normes.}

\begin{defi}{Normes équivalentes.}{}
    Soient $E$ un $\K$-ev et $\|\cdot\|_1$, $\|\cdot\|_2$ deux normes sur $E$.\\
    On dit que les deux normes sont \bf{équivalentes} ssi $\exists a,b > 0, ~ \forall x \in E, ~ a\|x\|_1\leq\|x\|_2\leq b\|x\|_1$.\\
    \bf{Remarque.} Cela définit une relation d'équivalence.
\end{defi}

\vspace*{-0.7cm}

\begin{ex}{}{}
    Vérifier que les trois normes usuelles sur $\R^n$ sont équivalentes.
    \tcblower
    Soit $x=(x_1,...,x_n)\in\R^n$. On a $\|x\|_1=\sum_{i=1}^n|x_i|$, $\|x\|_2=\sqrt{\sum_{i=1}^nx_i^2}$ et $\|x\|_{\infty}=\max_i(|x_i|)$.\\
    On a $\forall x \in \R^n, ~ \|x\|_1\leq n\|x\|_\infty$ et $\|x\|_\infty \leq \|x\|_1$, donc $\|x\|_1\sim \|x\|_\infty$.\\
    On a $\forall x \in \R^n, ~ \|x\|_2\leq \sqrt(x)\|x\|_{\infty}$, et $\|x\|_{\infty}\leq\|x\|_2$, donc $\|x\|_2\sim \|x\|_\infty$. 
\end{ex}

\begin{prop}{}{}
    Soit $E$ un $\K$-ev.
    \begin{enumerate}
        \item Soit $A\subset E$. Le fait que $A$ soit borné ne dépend pas de la norme équivalente choisie.
        \item Soit $(u_n)_{n\in\N}\in E^\N$. Le fait que $(u_n)_{n\in\N}$ soit bornée ne dépend pas de la norme équivalente choisie.
        \item Soit $(u_n)_{n\in\N}\in E^\N$. Le fait que $(u_n)_{n\in\N}$ converge et sa limite ne dépendent pas de la norme éq. choisie.
    \end{enumerate}
    \tcblower
    \boxed{1.} Supposons $A$ borné pour $\|\.\|_1$, et $\|\.\|_1\sim\|\.\|_2$.\\
    $\hra$ $\exists M \geq 0, ~ \forall x \in A, ~ \|x\|_1\leq M$ et $\exists b > 0, ~ \forall x \in E, ~ \|x\|_2 \leq b\|x\|_1$.\\
    Alors $\forall x \in E, ~ \|x\|_2 \leq bM$, donc $A$ est borné pour $\|\.\|_2$.\\
    \boxed{2.} En posant $A=\{u_n, ~ n\in\N\}$, on retrouve le cas précédent.\\
    \boxed{3.} Supposons $\|u_n-l\|_1\to0$ et $\|\.\|_1\sim\|\.\|_2$.\\
    $\hra$ $\exists a,b>0, ~ \forall x \in E, ~ a\|x\|_1\leq\|x\|_2\leq b\|x\|_1$.\\
    Alors $a\|u_n - l\|_1 \leq \|u_n-l\|_2 \leq b\|u_n-l\|_1$ et par encadrement, $\|u_n-l\|_2\to0$.
\end{prop}

\vspace*{-0.7cm}

\begin{ex}{}{}
    Soit $E=\R[X]$, muni des normes $\|P\|_1=\sum|a_k|$ et $\|P\|_\infty=\max |a_k|$.\\
    Montrer que ces normes ne sont pas équivalentes.
    \tcblower
    Posons $P_n=1+...+X^n\in E$. On a $\|P_n\|_\infty=1$ et $\|P\|_1=n+1\to+\infty$ donc $(P_n)$ non bornée dans $(E,\|\cdot\|_1)$.
\end{ex}

\vspace*{-0.7cm}

\begin{ex}{}{}
    Soit $E=\{(u_n)_{n\in\N} \in \R^\N, ~ \sum|u_n| \nt{ converge}\}$, muni de $\|(u_n)\|_1=\sum|u_n|$ et $\|(u_n)\|_{\infty}=\sup|u_n|$.\\
    Montrer que ces normes ne sont pas équivalentes.
\end{ex}

\vspace*{-0.7cm}

\begin{exi}{}{}
    Soit $E=\cursive{C}^1([0,1],\R)$ muni de $\|\cdot\|_\infty$ et de $\|f\|=|f(0)|+\sup|f'|$.
    \tcblower
    $\hra$ Pour $f\in\cursive{C}^1([0,1],\R)$, $f'$ est continue donc bornée sur $[0,1]$, donc $\|f\|$ est bien définie.\\
    Par TFA :
    \begin{equation*}
        \forall x \in [0,1], ~ f(x) = f(0) + \int_0^xf'(t)\dt ~\nt{donc}~ |f(x)|\leq|f(0)|+\int_0^x|f'(t)|\dt\leq|f(0)|+\int_0^1|f'(t)|\dt.
    \end{equation*}
    Donc $|f(x)|\leq\|f\|$, donc $\|f\|_{\infty} \leq \|f\|$.\\
    $\hra$ Posons $f_n(x)=\sin(\frac{n\pi}{2}x)\in E$ pour $x\in[0,1]$ et $n\in\N$.\\
    $\bullet$ $\forall n \in \N, ~ \|f_n\|_{\infty}=1$.\\
    $\bullet$ $\forall x \in \R, ~ f_n'(x) = \frac{n\pi}{2}\cos(\frac{n\pi}{2}x)$ donc $\|f_n\|=|f_n(0)|+n\frac{\pi}{2}\to+\infty$.\\
    Les deux normes ne sont pas équivalentes.
\end{exi}

\section{Topologie.}

\subsection{Ouverts, fermés.}

\begin{defi}{Voisinage.}{}
    Soit $(E,\|\cdot\|)$ evn et $a\in E$, $V\subset E$ est un \bf{voisinage} de $a$ ssi $\exists r > 0, ~ B(a,r)\subset V$.\\
    On notera $\m{V}(a)$ l'ensemble des voisinages de $a$.\\
    \bf{Remarque.} On a $E\in\m{V}(a)$ donc $\m{V}(a)\neq\0$.
\end{defi}

\vspace*{-0.8cm}

\begin{defi}{Ouvert.}{}
    Soit $(E,\|\cdot\|)$ un evn, alors $O\subset E$ est un \bf{ouvert} de $E$ ssi $\forall x \in O, ~ O\in\m{V}(x)$.
\end{defi}

\vspace*{-0.8cm}

\begin{defi}{Fermé.}{}
    Soit $(E,\|\cdot\|)$ un evn, alors $F\subset E$ est un \bf{fermé} ssi $E\setminus F$ est un ouvert.
\end{defi}

\vspace*{-0.8cm}

\begin{prop}{}{}
    Les boules ouvertes sont des ouverts, les boules fermées et les singletons sont des fermés.
    \tcblower
    \boxed{1.} Soit $a\in E$, $r>0$ et $x\in B(a,r)$.\\
    $\hra$ $\|x-a\|<r$, donc on pose $r'=r-\|x-a\|>0$.\\
    $\bullet$ Montrons que $B(x,r')\subset B(a,r)$.\\
    $\hra$ Soit $y\in B(x,r')$, on a $\|y-a\|\leq \|y-x\|+\|x-a\| < r' + \|x-a\| < r$, donc $y\in B(a,r)$.\\
    Ainsi, $B(a,r)$ est un ouvert par définition.\\
    \boxed{2.} Soit $a\in E$, $r>0$ et $x\in E\setminus B'(a,r)$.\\
    $\hra$ $\|x-a\|>r$, donc on pose $r'=\|x-a\|-r>0$.\\
    $\bullet$ Montrons que $B(x,r)\subset E\setminus B'(a,r)$.\\
    $\hra$ Soit $y\in B(x,r)$, on a $\|y-a\| \geq |\|x-a\| - \|y-x\|| > \|x-a\| - r' = \|x-a\| - (\|x-a\| - r) = r$.\\
    Ainsi, $B'(a,r)$ est un fermé par définition.\\
    \boxed{3.} Soit $a\in E$ et $x\in E\setminus\{a\}$.\\
    $\hra$ $\|x-a\|>0$, donc on pose $r=\|x-a\|$.\\
    Soit $y\in B(x,r)$, on a $y\neq a$, car sinon $\|x-a\|<r$, ce qui est absurde, donc $y\in B(x,r)$.\\
    Ainsi, $\{a\}$ est un fermé par définition.
\end{prop}

\vspace*{-0.8cm}

\begin{prop}{Réunions et intersections. $\star$}{}
    \begin{enumerate}
        \item Toute réunion d'ouverts est un ouvert.
        \item Toute intersection finie d'ouverts est un ouvert.
        \item Toute réunion finie de fermés est un fermé.
        \item Toute intersection de fermés est un fermé.
    \end{enumerate}
    \tcblower
    \boxed{1.} Soit $(O_i)_{i\in I}$ une famille d'ouverts, on pose $O$ leur réunion.\\
    $\hra$ Pour $x\in O$, il existe $i\in I$ tel que $x\in O_i$, et $O_i$ est un ouvert.\\
    Donc il existe $r_i > 0$ tel que $B(x,r_i)\subset O_i$, donc $B(x,r_i)\subset O$ : $O$ est un ouvert.\\
    \boxed{2.} Soit $(O_i)_{i\leq n}$ un famille finie d'ouverts, et $O$ leur intersection.\\
    $\hra$ Pour $x\in O, ~ \forall i \in \lb1,n\rb, ~ x \in O_i$ : $\exists r_i > 0, ~ B(x,r_i)\subset O_i$.\\
    Donc il existe $r=\min(r_1,...,r_n)$ tel que $\forall i\lb1,n\rb, ~ B(x,r)\subset O_i$, donc $B(x,r)\subset O$ et $O$ est un ouvert.\\
    \boxed{3.} Soit $(F_i)_{i \leq n}$ une famille finie de fermés, et $F$ leur réunion.\\
    $\hra$ $\forall i \in \lb1,n\rb, ~ \ov{F_i}$ est un ouvert, et $\ov{\bigcup F_i}=\bigcap\ov{F_i}$, intersection finie d'ouverts.\\
    \boxed{4.} idem au cas 2.
\end{prop}

\begin{ex}{}{}
    Les intervalles ouverts de $\R$ sont des ouverts.
\end{ex}

\vspace*{-0.8cm}

\begin{prop}{}{}
    Soit $E$ un evn, et $A\subset E$ fini. Alors $A$ est fermé.
    \tcblower
    Notons $n=|A|$, et $A=\{a_1,...,a_n\}=\bigcup \{a_i\}$; A est réunion finie de fermés.
\end{prop}

\vspace*{-0.8cm}

\begin{ex}{}{}
    Montrer que $\Z$ est un fermé.
    \tcblower
    On a $\ov{\Z}=\bigcup_{n\in\Z}]n,n+1[$ réunion d'ouverts donc ouvert, puis $\Z$ est fermé.
\end{ex}

\vspace*{-0.8cm}

\begin{prop}{Ouvert dans un espace produit.}{}
    Soit $(E_i,\|\cdot\|_i)_i$ des evn et $E=E_1\times...\times E_n$ l'evn produit.\\
    On rappelle que par définition, la norme sur $E$ est $\|x\|_\infty=\max\|x_i\|_i$ où $x=(x_1,...,x_n)$.\\
    Soit $O_i$ ouvert de $E_i$, alors $O=O_1\times...\times O_n$ est un ouvert de $E$.\\
    Un produit fini d'ouverts est un ouvert de l'espace produit.
    \tcblower
    Soit $(x_1,...,x_n)\in O$, $\forall i \in \lb1,n\rb, ~ O_i\in \m{V}(x_i)$ : $\exists r_i > 0, ~ \forall y \in E_i, ~ \|y-x_i\|<r_i \ra y \in O_i$.\\
    Soit $r=\min(r_1,...,r_n)$, et $y=(y_1,...,y_n)\in B(x,r)$. On a $\forall i \in \lb1,n\rb, ~ \|y_i-x_i\|_i < r$ donc $y_i \in O_i$ donc $y \in O$.
\end{prop}

\vspace*{-0.8cm}

\begin{prop}{}{}
    Soit $E$ un $\K$-ev, $\|\cdot\|_1$ et $\|\cdot\|_2$ deux normes sur $E$.
    \begin{enumerate}
        \item On suppose qu'il existe $a>0, ~ \forall x \in E, ~ \|x\|_1\leq a\|x\|_2$.\\
        Soit $O\subset E$, alors $O$ ouvert pour $\|\cdot\|_1 \ra O$ ouvert pour $\|\cdot\|_2$.
        \item On suppose $\|\cdot\|_1\sim\|\cdot\|_2$. Soient $O\subset E$ et $F\subset E$.\\
        Alors $O$ ouvert pour $\|\cdot\|_1\iff O$ ouvert pour $\|\cdot\|_2$ et $F$ fermé pour $\|\cdot\|_1\iff F$ fermé pour $\|\cdot\|_2$.
    \end{enumerate}
    \tcblower
    \boxed{1.} Soit $O\subset E$ ouvert dans $(E,\|\cdot\|_1)$.\\
    Soit $x\in O$. Par hypothèse, $\exists r_1 > 0, ~ \forall y \in E, ~ \|y-x\|_1<r_1 \ra y \in O$.\\
    On pose $r_2=\frac{r_1}{a}>0$, alors $\forall y \in E, ~ \|y-x\|_2<r_2 \ra \|y-x\|_1\leq a\|y-x\|_2<ar_2=r_1\ra y\in O$.\\
    \boxed{2.} Par hypothèse, $\exists a > 0, ~ \exists b > 0, ~ \forall x \in E, ~ a\|x\|_1\leq\|x_2\|\leq b\|x\|_1$.\\
    Donc $O$ est ouvert pour $\|\cdot\|_1 \ra O$ ouvert pour $\|\cdot\|_2$, puis $O$ est ouvert pour $\|\cdot\|_2 \ra O$ ouvert pour $\|\cdot\|_1$.\\
    Enfin, $F$ fermé pour $\|\cdot\|_1 \iff \leftindex^CF$ ouvert pour $\|\cdot\|_1 \iff \leftindex^CF$ ouvert pour $\|\cdot\|_2\iff F$ fermé pour $\|\cdot\|_2$.
\end{prop}

\vspace*{-0.8cm}

\begin{defi}{Ouverts relatifs.}{}
    Soit $(E,\|\cdot\|)$ un evn et $A\subset E$.\\
    Un \bf{ouvert relatif} de $A$ est un ensemble de la forme $O\cap A$ avec $O$ ouvert de $E$. 
\end{defi}

\vspace*{-0.8cm}

\begin{ex}{}{}
    Soit $A=[1,2]$. $O'=[1,\frac{3}{2}[$ est un ouvert relatif de $A$, en effet : $O'=~]0,\frac{3}{2}[~\cap~[1,2]$.
\end{ex}

\subsection{Intérieur, adhérence.}

À partir de maintenant, on note le complémentaire d'un ensemble $A$ par $\leftindex^C A$.

\begin{defi}{}{}
    Soit $(E,\|\cdot\|)$ un evn et $A\subset E$.\\
    $\hra$ L'\bf{intérieur} de $A$, noté ${\acc{A}}$, est le plus grand ouvert inclus dans $A$.\\
    $\hra$ L'\bf{adhérence} de $A$, notée $\ov{A}$, est le plus petit fermé contenant $A$.
    \begin{equation*}
        {\acc{A}}\subset A \subset \ov{A}
    \end{equation*}
\end{defi}

\vspace*{-0.8cm}

\begin{prop}{}{67}
    Soit $(E,\|\cdot\|)$ un evn et $A\subset E$.
    \begin{equation*}
        A \nt{ ouvert} \iff A = {\acc{A}}; \quad A \nt{ fermé} \iff A = \ov{A}.
    \end{equation*}
    \tcblower
    \boxed{\la} Supposons $A={\acc{A}}$, alors ${\acc{A}}$ est ouvert donc $A$ l'est aussi.\\
    \boxed{\ra} Supposons $A$ ouvert, alors $A={\acc{A}}$, le plus grand ouvert inclus dans $A$.
\end{prop}

\vspace*{-0.8cm}

\begin{defi}{Frontière.}{}
    Soit $(E,\|\cdot\|)$ evn et $A\subset E$. La \bf{frontière} de $A$, notée $\Fr(A)$ est l'ensemble $\Fr(A)=\ov{A}\setminus{\acc{A}}$.\\
    \bf{Remarque.} La frontière de $A$ est toujours fermée, comme intersection de fermés.
\end{defi}

\vspace*{-0.8cm}

\begin{thm}{}{67}
    Soit $(E,\|\cdot\|)$ et $A,B\in\P(E)$.
    \begin{enumerate}
        \item $A\subset B \ra {\acc{A}}\subset\accentset{\circ}{B}$.
        \item $A\subset B \ra \ov{A}\subset\ov{B}$.
    \end{enumerate}
    \tcblower
    \boxed{1.} On a ${\acc{A}}\subset A \subset B$ donc ${\acc{A}}$ ouvert de $B$, donc ${\acc{A}}\subset\accentset{\circ}{B}$.\\
    \boxed{2.} On a $A\subset B \subset \ov{B}$ donc $\ov{B}$ est un fermé contenant $A$ donc $\ov{A}\subset\ov{B}$.
\end{thm}

\vspace*{-0.8cm}

\renewcommand*{\A}{\accentset{\circ}{A}}
\renewcommand*{\c}{\leftindex^C}

\begin{lemme}{}{}
    Soit $(E,\|\cdot\|)$ un evn. Soit $A\subset E$.
    \begin{enumerate}
        \item $\leftindex^C{\ov{A}}=\accentset{\circ}{\leftindex^CA}$.
        \item $\leftindex^C{\acc{A}}=\ov{\leftindex^CA}$.
    \end{enumerate} 
    \tcblower
    $\hra$ $A\subset\ov{A}$ donc $\c{\ov{A}}\subset\c{A}$ donc $\c{\ov{A}}\subset\accentset{\circ}{\c{A}}$.\\
    $\hra$ $\accentset{\circ}{\c{A}}\subset\c{A}$ donc $A \subset \c{\left(\accentset{\circ}{\c{A}}\right)}$ donc $\ov{A}\subset \c{\left(\accentset{\circ}{\c{A}}\right)}$ donc $\acc{\c{A}} \subset \c{\ov{A}}$.\\
    Par double-inclusions : $\c{\ov{A}}=\acc{\c{A}}$.\\
    On utilise la proposition \ref{prop:67} sans détailler, et $B$ fermé (resp. ouvert) de $A$ implique $\ov{A}\subset B$ (resp. $B \subset \A$).\\
    $\hra$ $\A \subset A$ donc $\c{A}\subset\c{\A}$ donc $\ov{\c A} \subset \c{\A}$.\\
    $\hra$ $\c{A}\subset\ov{\c{A}}$ donc $\c{\ov{\c{A}}}\subset A$ donc $\c{\ov{\c{A}}}\subset\A$ puis $\c{\A}\subset\ov{\c{A}}$.\\
    Par double-inclusions : $\c{\A}=\ov{\c{A}}$.
\end{lemme}

\pagebreak

\begin{thm}{Caractérisation des éléments de l'intérieur et de l'adhérence. $\star$}{caradh}
    Soit $(E,\|\cdot\|)$ un evn et $A\subset E$.
    \begin{enumerate}
        \item $x\in{\acc{A}} \iff \exists V \in \m{V}(x), ~ V \subset A \iff \exists r>0, ~ B(x,r)\subset A$.\\
        Les éléments de l'intérieur de $A$ sont les éléments qui sont à l'intérieur.
        \item $x\in\ov{A}\iff \forall V \in \m{V}(x), ~ V \cap A \neq 0 \iff \forall r > 0, ~ B(x,r)\cap A \neq 0$.\\
        Les éléments de l'adhérence de $A$ sont les éléments adhérents à la partie.
    \end{enumerate}
    \tcblower
    \boxed{1.} La deuxième équivalence puis la première:\\
    \boxed{\ra} $\exists V \in \m{V}(x), ~ V\subset A$ puis $\exists r > 0, ~ B(x,r)\subset V$ donc $B(x,r)\subset A$.\\
    \boxed{\la} $\exists r > 0, ~ B(x,r)\subset A$ et $B(x,r)$ est un voisinage de $x$.\\
    \boxed{\ra} Soit $x\in{\acc{A}}$ ouvert donc $\exists r>0, ~ B(x,r)\subset {\acc{A}}$, c'est un voisinage.\\
    \boxed{\la} $\exists r > 0, ~ B(x,r)\subset A$ donc $B(x,r)$ est un ouvert de $A$, donc $B(x,r)\subset{\acc{A}}$.\\
    \boxed{2.} On utilise la 1. et le passage au complémentaire.\\
    On a $x\notin\ov{A}\iff x\in\leftindex^C{\ov{A}}\iff x \in \accentset{\circ}{\leftindex^CA}\iff \exists r>0, ~ B(x,r)\subset\leftindex^CA$.\\
    Alors $x\in \ov{A} \iff \forall r > 0, ~ B(x,r) \cap A \neq 0$.
\end{thm}

\begin{thm}{Caractérisation séquentielle de l'adhérence. $\star$}{seqadh}
    Soit $(E,\|\cdot\|)$ un evn et $A\subset E$.
    \begin{equation*}
        x \in \ov{A} \iff \exists (a_n)_{n\in\N} \in A^\N, ~ \lim a_n = x.
    \end{equation*}
    \tcblower
    \boxed{\la} Soit $(a_n)\in A^\N$ telle que $\lim a_n = x$. Soit $r>0, ~ \exists n_0 \in \N, ~ \forall n \geq n_0, ~ \|a_n-x\| < r$.\\
    Alors $a_N\in B(x,r)\cap A\neq\0$, donc $x\in\ov{A}$.\\
    \boxed{\ra} Soit $x\in\ov{A}$. On cherche à construire une suite $(a_n)$ qui convient.\\
    Par hypothèse, $\forall r > 0, ~ B(x,r)\cap A \neq 0$, donc pour $n\in\N^*, ~ B(x,\frac{1}{n})\cap A\neq\0$.\\
    Ainsi, $\forall n \in \N^*, ~ \exists a_n \in B(x,\frac{1}{n})\cap A$, on a construit la suite $(a_n)$, et par encadrement, $a_n\to x$.
\end{thm}

\vspace*{-0.4cm}

\begin{corr}{Caractérisation séquentielle d'un fermé.}{}
    Soit $(E,\|\cdot\|)$ un evn et $F\subset E$.
    \begin{center}
        $F$ est fermé $\iff$ toute limite de suite convergente d'éléments de $F$ est dans $F$.
    \end{center}
    \tcblower
    On a toujours $F\subset\ov{F}$ donc $F=\ov{F} \iff \ov{F} \subset F \iff $ l'énoncé.
\end{corr}

\vspace*{-0.4cm}

\begin{ex}{}{}
    $\hra$ $[1,2]$ est un fermé de $\R$ puisque pour $(u_n)\in[1,2]^\N$ convergente vers $l$ :
    \begin{equation*}
        \forall n \in \N, ~ 1 \leq u_n \leq 2 \ra 1 \leq l \leq 2.
    \end{equation*}
    $\hra$ $\Z$ est un fermé de $\R$ car pour $(u_n)\in\Z^\N$, convergente vers $l$, on a :
    \begin{equation*}
        \exists n_0 \in \N, ~ \forall n \geq n_0, ~ |u_n-l| <\frac{1}{2}, ~ \nt{donc} ~ |u_n - u_{n_0}| < 1.
    \end{equation*}
    Donc $\forall n \geq n_0, ~ u_n = u_{n_0}$, donc $l=u_{n_0}\in\Z$.
\end{ex}

\vspace*{-0.4cm}

\begin{ex}{}{}
    Soit $(E,\|\cdot\|)$ un evn et $F$ un ssev de $E$. Montrer que $\ov{F}$ est un ssev de $E$.
    \tcblower
    On a $\ov{F}\neq\0$ car $0_F\in F\subset \ov{F}$. Soient $x,y\in \ov{F}$ et $\l\in\K$.\\
    On a $\exists (a_n)\in F^\N, ~ x = \lim a_n$ et $\exists (b_n)\in F^\N, ~ y=\lim b_n$.\\
    On a $(a_n+\l b_n)\to x + \l y$ par théorème, or $a_n+\l b_n \in F$ donc $x+\l y \in \ov{F}$.
\end{ex}

\vspace*{-0.8cm}

\begin{defi}{Distance à une partie.}{}
    Soit $(E,\|\cdot\|)$ un evn, et $A\subset E$ non-vide. Soit $x\in E$.\\
    La \bf{distance} de $x$ à $A$ est définie par $d(x,A)=\inf\{\|x-y\|, ~ y \in A\}$.
    \tcblower
    Notons $D=\{\|x-y\|, ~ y \in A\}$, on a $D\subset\R_+$ et $A$ non-vide donc $y_0\in A$ puis $\|x-y_0\|\in D$ donc $D\neq\0$.\\
    On a alors $D$ une partie non-vide de $\R$ et minorée par 0, donc l'inf existe par théorème. 
\end{defi}

\vspace*{-0.8cm}

\begin{ex}{}{}
    Soit $(E,\|\cdot\|)$ un evn et $A\subset E$ non-vide. Montrer que :
    \begin{equation*}
        x \in \ov{A} \iff d(x,A) = 0.
    \end{equation*}
    \tcblower
    \boxed{\ra} On a $\exists (a_n) \in A^\N, ~ a_n \to x$. Or $\forall n \in \N, ~ 0 \leq d(x,A)\leq \|x-a_n\|$, puis par encadrement, $d(x,A)=0$.\\
    \boxed{\la} On a $d(x,A)=0$ donc $\exists (\a_n)\in \{\|x-y\|\}^\N$ telle que $\lim \a_n = \inf \{\|x-y\|, ~ y \in A\}=0$.\\
    Donc $\forall n \in \N, ~ \a_n\in D$, donc $\exists y_n \in A, ~ \a_n = \|x-y_n\|$.\\
    On a donc $\lim \|x-y_n\|=0$ donc $y_n\to x$ donc $x\in\ov{A}$.
\end{ex}

\vspace*{-0.8cm}

\begin{ex}{}{}
    Soit $(E,\|\cdot\|)$ un evn, $a\in E$ et $r>0$.
    \begin{enumerate}
        \item $\accentset{\circ}{B'(a,r)}=B(a,r)$ et $\ov{B'(a,r)}=B'(a,r)$.
        \item $\ov{B(a,r)}=B'(a,r)$ et $\accentset{\circ}{B(a,r)}=B(a,r)$.
        \item $\Fr(B(a,r))=\Fr(B'(a,r))=S(a,r)$.
    \end{enumerate}
    \tcblower
    \boxed{1.} On a déjà $B'(a,r)$ est un fermé en tant que boule fermée, donc $\ov{B'(a,r)}=B'(a,r)$.\\ 
    \boxed{2.} On a déjà $B(a,r)$ est un ouvert en tant que boule ouverte, donc $\accentset{\circ}{B(a,r)}=B(a,r)$.\\
    \boxed{\subset} On a $B'(a,r)$ est un fermé, qui contient $B(a,r)$, donc $\ov{B(a,r)}\subset B'(a,r)$.\\
    \boxed{\supset} Soit $x\in B'(a,r)$ : $\|x-a\|\leq r$. Si $x\in B(a,r)$, alors $x\in\ov{B(a,r)}$.\\
    Sinon, $x\in S(a,r)$ donc $\|x-a\|=r$. On cherche une suite $(a_n)_{n\in\N}$ de $B(a,r)$ convergente vers $x$.\\
    $\bullet$ Pour tout $n\in\N^*$, posons $a_n = \frac{1}{n}a + \left( 1 - \frac{1}{n} \right)x$.\\
    $\hookrightarrow$ $\forall n \geq 1, ~ a_n - a = \left( \frac{1}{n} - 1 \right)(a-x)$ puis $\|a_n - a\|=\left( 1 - \frac{1}{n} \right)\|x-a\|<r$ donc $a_n \in B(a,r)$.\\
    $\hookrightarrow$ $\forall n \geq 1, ~ a_n - x = \frac{1}{n}(a-x)$ puis $\|a_n-x\|=\frac{1}{n}\|a-x\|=\frac{r}{n}\to0$, donc $a_n\to x$.\\
    On a bien $x\in \ov{B(a,r)}$ par caractérisation.\\
    \boxed{3.} Évident car $\Fr(A)=\ov{A}\setminus{\acc{A}}=S(a,r)$, et on utilise 1. et 2.
\end{ex}

\vspace*{-0.7cm}

\begin{defi}{Densité.}{}
    Soit $(E,\|\cdot\|)$ un evn et $A\subset E$. On dit que $A$ est dense ssi l'une des propositions est vraie :
    \begin{enumerate}
        \item $\ov{A}=E$.
        \item $\forall x \in E, ~ \exists (a_n) \in A^\N, ~ \lim a_n = x$.
        \item $\forall O$ ouvert non-vide de $E$, $O\cap A \neq \0$.
        \item $\forall x \in E, ~ \forall r > 0, ~ B(x,r) \cap A \neq \0$
    \end{enumerate}
    Ces équivalences sont justifiées par les théorèmes \ref{thm:caradh} et \ref{thm:seqadh}.
\end{defi}

\vspace*{-0.7cm}

\begin{prop}{}{}
    $\Q$ et $\R\setminus\Q$ sont denses dans $\R$.\\
    \bf{Remarque.} Tout intervalle ouvert non vide de $\R$ contient une infinité de rationnels.
    \tcblower
    Soient $a,b\in\R$ tels que $a<b$. Montrons que $]a,b[\cap\Q\neq\0$.\\
    \bf{Analyse.} Soient $p,q\in\Z\times\N^*$ tels que $p\land q = 1$ et $a<\frac{p}{q}<b$, alors $aq<p<bq$.\\
    Il suffit de choisir $q$ tel que $bq-aq>1$ donc $q>\frac{1}{b-a}$.\\
    \bf{Synthèse.} Choisissons un entier $q$ tel que $q>\frac{1}{b-a}$, possible car $b-a>0$.\\
    Puis on choisit un entier $p$ entre $aq$ et $bq$ : $\lf aq \rf + 1$.\\
    On a bien $\lf aq \rf \leq aq < p$ donc $aq < p < aq + 1$ donc $a<\frac{p}{q}<a+\frac{1}{q}$ donc $a<\frac{p}{q}<b$.\\
    $\bullet$ Cherchons un irrationnel : $\sqrt{2}\in\R\setminus\Q$ puis $\exists x \in \Q, ~ \frac{a}{\sqrt{2}} < x < \frac{b}{\sqrt{2}}$, puis $a<x\sqrt{2}<b$.
\end{prop}

\vspace*{-0.7cm}

\subsection{Compacité.}

\begin{defi}{Compacité.}{}
    Soit $(E,\|\cdot\|)$ un evn et $K\subset E$.
    \begin{center}
        $K$ est \bf{compact} ssi de toute suite d'éléments de $K$, on peut extraire une suite convergente dans $K$.
    \end{center}
    \bf{Remarque.} L'ensemble vide et les singletons sont compacts.
\end{defi}

\begin{ex}{}{}
    Les segments de $\R$ sont compacts.
    \tcblower
    Soit $[a,b]$ un segment de $\R$ et $(u_n)_{n\in\N}\in[a,b]^\N$. Alors $\forall n \in \N, ~ a \leq u_n \leq b$ :  la suite est bornée.\\
    Ainsi, par théorème de Bolzano-Weierstraß, on peut extraire de $(u_n)_{n\in\N}$ une suite convergente de limite $l\in[a,b]$.  
\end{ex}

\begin{prop}{}{}
    Soit $(E,\|\cdot\|)$ un evn. Toute partie finie de $E$ est compacte.
    \tcblower
    Soit $K=\{a_1,...,a_p\}$ une partie finie de $E$, de cardinal $p$.\\
    Soit $(u_n)_{n\in\N}\in K^\N$ et $\forall i \in \lb1,p\rb, ~ F_i=\{n\in\N, ~ u_n = a_i\}$. On a $\N=F_1\sqcup...\sqcup F_p$.\\
    On a $\exists i_0 \in \lb1,p\rb, ~ F_{i_0}$ de cardinal infini, car sinon $\N$ est fini.\\
    Ainsi, on pose $\phi(0)=\min F_{i_0}$, ... puis $\phi(k+1)=\min F_{i_0} \setminus \{\phi(0),...,\phi(k)\}$ strictement croissante.\\
    Alors $(u_{\phi(n)})_{n\in\N}$ est constante égale à $a_{i_0}$, elle converge vers $a_{i_0}\in K$.
\end{prop}

\begin{prop}{}{}
    Soit $(E,\|\cdot\|)$ un evn. Toute partie compacte est fermée et bornée.\\
    \bf{Remarque.} La réciproque est fausse en dimension infinie.
    \tcblower
    Soit $K$ une partie compacte de $E$. Montrons que $K$ est fermée par caractérisation séquentielle.\\
    Soit $(u_n)_{n\in\N}\in K^\N$ convergente de limite $l$. Il existe $\phi:\N\to\N$ strictement croissante et $l'\in K$ tq $u_{\phi(n)}\to l'$.\\
    Or $u_{\phi(n)}\to l$ donc $l=l'$ par unicité de la limite, puis $l\in K$.\\
    Par contraposée, supposons $K$ non borné, montrons que $K$ n'est pas compact. Alors $\forall M > 0, ~ \exists x \in K, ~ \|x\|>M$.\\
    Alors $\forall n \in \N^*, ~ \exists x_n \in K, ~ \|x_n\|>n$. On a construit la suite $(x_n)$.\\
    Alors $\forall \phi : \N \to \N$ strictement croissante, $\forall n \in \N^*, ~ \|x_{\phi(n)}\|>\phi(n)\geq n$, donc $\|x_{\phi(n)} \| \to +\infty$.\\
    Aucune suite extraite de $(x_n)_{n\in\N}$ n'est bornée, donc aucune suite extraite ne converge, donc $K$ n'est pas compacte.
\end{prop}

\vspace*{-0.8cm}

\begin{ex}{}{}
    Soit $\|\cdot\|$ une norme sur $\R[X]$ définie par $\|P\|=\max(|a_k|, ~ k \in \N)$.\\
    Soit $S=\{P \in E, ~ \|P\|=1\}$ la sphère unité, $S$ est fermée et bornée.\\
    Vérifions que $S$ est non-compacte.
    \tcblower
    Pour tout $n\in\N$, posons $P_n=1+X+...+X^n$, on a $\|P_n\|=1$ donc $(P_n)_n\in S^\N$.\\
    On a $n \neq m$ (par exemple $n<m$) implique $\|P_m - P_n\|=1$, donc $\forall \phi, ~ n \neq m \ra \|P_{\phi(m)} - P_{\phi(n)}\|=1$.\\
    Supposons par l'absurde qu'il existe $l\in E, ~ P_{\phi(n)}\to l$. Alors $\exists N \in \N, ~ \forall n \in \N, ~ n\geq N \ra \|P_{\phi(n)}-l\|<\frac{1}{4}$.\\
    Donc $n,m>N\ra \|P_{\phi(m)-P_{phi(n)}}\|\leq \|P_{\phi(n)} - l\| + \|P_{\phi(m)} - l\| < \frac{1}{2} \ra 1 < \frac{1}{2}$, ...
\end{ex}

\vspace*{-0.8cm}

\begin{prop}{$\star$}{}
    Soit $(E,\|\cdot\|)$ un evn, $K$ un compact et $F$ fermé tel que $F\subset K$, alors $F$ est compact.
    \tcblower
    Soit $(u_n)\in F^\N$, alors $\forall n \in \N, ~ u_n \in K$. Or $K$ compact, donc $\exists \phi, \exists l \in K, ~ u_{\phi(n)}\to l$.\\
    Puis $l=\lim u_{\phi(n)}\in F$ car $F$ est fermé et $u_{\phi(n)}\in F$, donc $F$ compact.
\end{prop}

\vspace*{-0.8cm}

\begin{ex}{}{}
    Soit $(E,\|\cdot\|)$ un evn.
    \begin{enumerate}
        \item Soit $K\subset E$ un compact. $\leftindex^CK$ est-il compact ? 
        \item Soient $K_1,K_2$ deux compacts. $K_1\cap K_2$ est-il compact ?
        \item Soit $(K_i)_{i\in I}$ une famille de compacts. $\bigcap K_i$ est-il compact ?
        \item Soient $K_1,K_2$ deux compacts. $K_1\cup K_2$ est-il compact ?
        \item Soit $(K_i)_{i\in I}$ une famille de compacts. $\bigcup K_i$ est-il compact ?
    \end{enumerate}
    \tcblower
    \boxed{1.} Puisque $K$ est borné, $\leftindex^CK$ ne l'est pas, sauf si $E=\{0\}$ car un ev n'est pas borné.\\
    \boxed{2.} On a $K_1\cap K_2 \subset K_1$, fermé comme intersection, on utilise la proposition précédente.\\
    \boxed{3.} Si $I$ est vide, $\0$ est compact, sinon $\bigcap K_i$ est un fermé, inclus dans l'un des $K_i$, c'est un compact.\\
    \boxed{4.} Soit $(u_n)\in(\K_1\cup K_2)^\N$, $F_1=\{n\in\N, ~ u_n\in K_1\}$ et $F_2=\{n\in\N, ~ u_n \in K_2\}$..\\
    On a $F_1\cup F_2 = \N$, donc $F_1$ ou $F_2$ est infini, par exemple $F_1$. On a alors $(u_{\phi(n)})\in K_1^\N$ extraite.\\
    Puis $K_1$ compact donc on peut extraire une suite $(u_{\phi\circ\psi(n)})$ convergente dans $K_1$, donc convergente dans $K_1\cup K_2$.\\
    \boxed{5.} Non, par exemple : $\bigcup_{n\in\Z}[n,n+1]=\R$ pas un compact car non-borné.\\
    \bf{Remarque.} On peut généraliser 2. et 4. par récurrence.
\end{ex}

\begin{prop}{Produit de compacts. $\star$}{}
    Soit $(E_i, \|\cdot\|_i)_{i\leq n}$ des evn, $K_i$ un compact de $E_i$ pour tout $i$, et $E=E_1\times...\times E_n$ l'evn produit.
    \begin{center}
        \fbox{$K=K_1\times...\times K_n$ est un compact de $E$.}
    \end{center}
    \tcblower
    Par récurrence sur $n$.\\
    \bf{Initialisation.} Soient $E=E_1\times E_2$, $K_1$ compact de $E_1$, $K_2$ compact de $E_2$, $K=K_1\times K_2$.\\
    Soit $(a_n)_{n\in\N}\in K^\N$. Alors $\forall n \in \N, \exists (x_n, y_n) \in K_1 \times K_2, ~ a_n = (x_n, y_n)$.\\
    On effectue un procédé d'extraction diagonale :\\
    --- $K_1$ est compact : $\exists \phi, ~ \exists l \in K_1, ~ x_{\phi(n)}\to l_1$.\\
    --- $K_2$ est compact, et $(y_{\phi(n)})\in K_2^\N$, donc $\exists \psi, ~ \exists l_2 \in K_2, ~ y_{\phi\circ\psi(n)}\to l_2$.\\
    Puis $x_{\phi\circ\psi(n)}\to l_1$ en tant que suite extraite de $(x_{\phi(n)})$, donc $a_{\phi\circ\psi(n)}\to(l_1,l_2)$, puis $K_1\times K_2$ compact.\\
    \bf{Hérédité.} Soit $n\in\N$, supposons que $\forall E_1,...,E_n$, $\forall K_1,...,K_n$, $K_1\times...\times K_n$ est compact de $E_1\times...\times E_n$.\\
    Soient $E_1,...,E_{n+1}$ des evn et $K_i$ compact de $E_i$ pour tout $i$. Alors :\\
    --- $K_1\times ... \times K_n$ est un compact de $E_1\times...\times E_n$ par hypothèse.\\
    --- $K_n+1$ est un compact de $E_{n+1}$.\\
    D'après la propriété au rang 2, $K_1\times...\times K_{n+1}$ est compact de $E_1\times...\times E_{n+1}$.\\
    On conclut par récurrence.
\end{prop}

\begin{prop}{}{}
    Soit $(E,\|\cdot\|)$ un evn. Soit $K\subset E$ un compact et $(u_n)_n\in K^\N$.\\
    Alors $(u_n)$ converge ssi $(u_n)_n$ admet une unique valeur d'adhérence.
    \tcblower
    \boxed{\ra} Si $(u_n)_n$ converge vers $l$, alors ses suites extraites aussi : $\Val((u_n))=\{l\}$.\\
    \boxed{\la} Supposons que $\Val(u_n)=\{l\}$. Par l'absurde, on suppose que $\exists \e > 0, ~ \forall N \in \N, ~ \exists n \geq N, ~ \|u_n - l\| > \e$.\\
    Équivalent à $\exists \e > 0, ~ \{n \in \N, ~ \|u_n - l\| > \e\}$ est infini puis $\exists\e>0, ~ \exists \phi , ~ \forall n \in \N, ~ \|u_{\phi(n)}-l\|>\e$.\\
    Comme $K$ est compact, $\exists \psi, ~ u_{\phi\circ\psi(n)}\to l' \in K$, donc $l'\neq l$, puis $l'\in \Val(u_n)$, absurde.
\end{prop}

\begin{ex}{}{}
    Soit $(E,\|\cdot\|)$ un evn et $(u_n)_{n\in\N}\in E^\N$ convergente vers $l\in E$. Montrer que $K=\{u_n, ~ n\in\N\}\cup\{l\}$ est un compact.
    \tcblower
    \bf{Idée.} Pour $(v_n)\in K^\N$, en extraire une suite convergente dans $K$. 
\end{ex}

\pagebreak

\section{Khôlles.}

\begin{exercice}{}{}
    Soit $E$ un evndf, $f\in\L(E)$ tel que $\forall x \in E, ~ \|f(x)\|\leq\|x\|$.\\
    Montrer que $E=\Ker(f-\id)\oplus\Im(f-\id)$.
    \tcblower
    Soit $x\in\Ker(f-\id)\cap\Im(f-\id)$, alors $\exists y \in E, ~ x = f(y)-y$, or $f(x)-x=0$ donc $x=f(x)$.\\
    Par récurrence sur $n\in\N^*$, on pose $\P_n$ : << $f^n(x+y) = (n+1)x-y$ >>.\\
    \bf{Initialisation.} $f(x+y)=f(x)+f(y)=f(y)+x=2x-y$. Donc $\P_1$ est vraie.\\
    \bf{Hérédité.} Soit $n\in\N$, on suppose $\P_n$.
    \begin{align*}
        f^{n+1}(x+y)&=f^{n+1}(x)+f^{n+1}(y)=x+f^n(f(y))\\
        &=x+f^n(x+y)=x+(n+1)x-y\\
        &=(n+2)x-y
    \end{align*}
    Par récurrence, $\forall n \in \N, ~ \P_n$.\\
    Soit $n\in\N$, $\|f^n(y)\|\leq\|y\|$ (récurrence facile) puis $\|f^{n-1}(x+y)\|\leq\|y\|$.\\
    Alors $\|nx+y\|\leq\|y\|$ et $n\|x\|\leq2\|y\|$. Puis $\|x\|\leq\frac{2}{n}\|y\|\to0$ donc $x=0$ par séparation.\\
    Alors $E=\rg(f-\id)+\dim\Ker(f-\id)$, puis par caractérisation, $E=\Ker(f-\id)\oplus\Im(f-\id)$.
\end{exercice}

\begin{exercice}{}{}
    Soit $E\subset\R^\N$ le $\R$-ev des suites réelles bornées. Pour $u\in E$, on définit :
    \begin{equation*}
        \|u\|=\sum_{n=0}^{+\infty}\frac{|u_n|}{2^n}
    \end{equation*}
    Vérifier que $\|\cdot\|$ définit une norme et montrer que $A=\{u\in E, ~ \forall n \in \N, ~ u_n \in [0,1]\}$ est compact.
    \tcblower
    La série est bien définie : pour $u\in E$, $|u_n|$ est bornée donc $\exists M > 0, ~ \frac{|u_n|}{2^n}\leq\frac{M}{2}$.\\
    Par comparaison (positifs) à une série convergente comme série géométrique, la série converge.\\
    \bf{Séparation.} Supposons $\|u\|=0$, alors $\sum_{n=0}^{+\infty}\frac{|u_n|}{2^n}=0$, donc $\forall n \in \N, ~ \frac{|u_n|}{2^n}=0$ puis $u_n=0$, donc $u_n=0_E$.\\
    \bf{Homogénéité, Inégalité triangulaire.} [...]\\
    \bf{Compact.} Soit $(u^k)\in A^\N$. Montrons que $A$ est compact.\\
    On a $(u_n^0)=(u_0^0,u_1^0,...)\in[0,1]^\N$ d'après Bolzano-Weierstraß : $\exists \phi_0$ strct. crs. $u^0_{\phi_0(n)}\to u^0\in\R$.\\
    On a $(u_{\phi_0(n)^1})=(u_{\phi_0(0)^1}, u_{\phi_0(1)^1}, ...)\in [0,1]^\N$, d'après Bolzano-Weierstraß : $\exists \phi_1$ strct. crs. $u_{\phi_0\circ\phi_1(n)}^1\to u^1\in\R$.\\
    (...) 
\end{exercice}

\pagebreak

\begin{exercice}{}{}
    Soit $f:\R\to\R$ de classe $\m{C}^0$. Soit $K$ un compact. Montrer que :
    \begin{equation*}
        f^{-1}(K) \nt{ est compact} \iff \lim_{x\to\pm\infty}|f(x)|=+\infty.
    \end{equation*}
    \tcblower
    \boxed{\ra} Par l'absurde, on suppose que $\lim_{x\to\pm\infty}|f(x)|\neq \infty$.\\
    Si $f(x)\to l \in \R_+$ en l'infini, alors : $\forall \e>0, ~ \exists M \geq 0, ~ \forall x \geq M, ~ |f(x)-l|\leq \e$.\\
    Donc $f(x)\in[l-\e,l+\e], ~ \forall x \geq M$, donc $f^{-1}([l-\e,l+\e])=[M,+\infty[$, absurde.\\
    \boxed{\la} Supposons que $\lim_{x\to\pm\infty}|f(x)|=+\infty$, et soit $K$ un compact non-vide.\\
    Si $\lim_{x\to\pm\infty}f(x)=+\infty$ (ou $-\infty$) spdg, alors, par continuité, il existe $x_0\in\R, ~ f(x_0)=\min_\R f$ (TBA).\\
    --- Si $\forall x \in K, ~ x < x_0$, alors $f^{-1}(K)=\0$, compact.\\
    --- Sinon, $\exists x \in K, ~ x \geq x_0$, soit $x_0=\min(K)\et x_1=\max(K)$, puis [...]. 
\end{exercice}

\vspace*{-0.7cm}

\begin{exercice}{}{}
    On munit $M_n(\R)$ du produit défini par $(A\mid B)=\Tr(A^\top B)$.
    \begin{enumerate}
        \item Montrer que la base canonique $(E_{i,j})_{1\leq i,j \leq n}$ de $M_n(\R)$ est orthonormée.
        \item Montrer que $S_n(\R)$ et $A_n(\R)$ sont supplémentaires orthogonaux.
        \item Montrer que $\forall A \in M_n(\R)$, $\Tr(A)\leq\sqrt{n}\sqrt{\Tr(A^\top A)}$ et préciser le cas d'égalité.
    \end{enumerate}
    \tcblower
    \boxed{1.} Soit $E_{i,j}$ avec $i,j\in\lb1,n\rb$. Alors $(E_{i,j}\mid E_{j,i})=\Tr(E_{j,i}E_{i,j})=\Tr(E_{j,j})=1$.\\
    Soient $E_{i,j},E_{k,l}$ avec $(i,j)\neq(k,l)$. Alors $(E_{i,j}\mid E_{k,l})=\Tr(\d_{i,k}E_{j,l})=\d_{i,k}\Tr(E_{j,l})$\\
    \boxed{2.} Soit $M\in M_n(\R)$, alors $M=\frac{M+M^\top}{2}+\frac{M-M^\top}{2}$.\\
    Soient $A\in S_n(\R)$ et $B\in A_n(\R)$. Alors $(A\mid B) = \Tr(A^\top B)=\Tr(AB)=\Tr(B^\top A^\top)=-\Tr(BA)$ donc $2\Tr(AB)=0$ puis $(A\mid B)=0$.\\
    \boxed{3.} On a $|(I_n\mid A)|\leq\sqrt{(I_n \mid I_n)}\sqrt{(A\mid A)}$ puis $\Tr(I_n^\top A) \leq \sqrt{\Tr(I_n)}\sqrt{\Tr(A^\top A)}$ puis $\Tr(A)\leq\sqrt{n}\sqrt{\Tr(A^\top A)}$.\\
    Supposons $A=\l I_n$, alors $\Tr(A)=\l n$ puis $\sqrt{n}\sqrt{\Tr(A^\top A)}=\sqrt{n}\sqrt{\Tr(\l^2 I_n)}=|\l|n$.
\end{exercice}

\begin{exercice}{}{}
    Déterminer l'adhérence de $A=\{(x,\cos\left( \frac{1}{x} \right)) \mid x\in\R, ~ x > 0\}$.
    \tcblower
    $\{(0,y), ~ y \in [-1,1]\}\cup\{(x,\cos(\frac{1}{x})), ~ x > 0\}$.
\end{exercice}

\begin{exercice}{}{}
    \small Soit $A$ une partie de $\R$. Un point $x\in\R$ est dit d'accumulation de $A$ si tout voisinage $x$ dans $\R$ privé de $x$ rencontre $A$. Montrer que si $A$ est bornée et n'admet qu'un seul point d'accumulation, alors il existe une bijection entre $\N$ et $A$.
\end{exercice}

\pagebreak

\begin{exercice}{(Enzo)}{}
    Soit $A\in M_n(\R)$ telle que $A^\top A=I_n$ et $\Ker(A-I_n)=\{0\}$.
    \begin{enumerate}
        \item Étudier la convergence de $u_p=\frac{1}{p+1}\left( I_n + ... + A^p \right)$.
        \item La suite $(A^p)_p$ est-elle convergente ?
    \end{enumerate}
    \tcblower
    \boxed{1.} On a $\|A\|=\sqrt{\Tr(A^\top A)}$.
    \begin{equation*}
        v_n := \frac{1}{p+1}\left( A - I_n \right)\sum_{k=0}^pA^k = \frac{1}{p+1}\left( A^{p+1} - I_n \right)
    \end{equation*}
    Or $\Tr((A^p)^\top A^p)=\Tr((A^{p-1})^\top A^\top A A^{p-1})=\Tr(I_n)=n$. Donc :
    \begin{equation*}
        \|v_n\|\leq\frac{1}{p+1}\|A^{p+1}\|+\|I_n\| \leq \frac{2\sqrt{n}}{p+1} \to_{p\to+\infty} 0
    \end{equation*}
    Ensuite, $\|(A-I)^{-1}v_p\|\leq\|(A-I_n)^{-1}\|\|v_p\|\leq\|(A-I_n)^{-1}\|\frac{2\sqrt{n}}{p+1}\to_{p\to+\infty}0$.\\
    Alors $\|u_p\|\to 0$ puis $u_p\to0$.\\
    \boxed{2.} [...]
\end{exercice}

\end{document}